\section{Cấu hình thí nghiệm}
\label{sec:setup}

\paragraph{Hạ tầng tính toán.}
Toàn bộ huấn luyện paper-faithful được thực hiện trên Kaggle Notebook GPU:
\begin{itemize}\itemsep0pt
    \item GPU: NVIDIA Tesla P100-PCIE-16GB (sm\_60, Pascal).
    \item Python 3.10, PyTorch 2.5.1+cu121 (downgrade từ default Kaggle 2.10 để hỗ trợ sm\_60).
    \item Mixed precision fp16 (AMP) cho phép batch 8 trong giới hạn 16GB VRAM.
\end{itemize}

\paragraph{Hyperparameters.}
Theo cấu hình paper sect.~5.2 (giữ nguyên với mọi biến thể để đảm bảo so sánh công bằng):
\begin{center}
\begin{tabular}{ll}
\toprule
Tổng epoch & 120 (Pha I: 40, Pha II: 80) \\
Optimizer  & 6 Adam riêng (Encoder, Decoder, BaseFuseLayer, DetailFuseLayer, AG-Base, AG-Detail) \\
Learning rate & $10^{-4}$, StepLR $\gamma = 0.5$ mỗi 20 epoch, floor $10^{-6}$ \\
Batch size & 8, AMP fp16 \\
Clip grad norm & 0.01 \\
$\alpha_1, \alpha_2, \alpha_3, \alpha_4$ & $1, 2, 5, 2$ (theo code release) \\
Seed & 42 \\
\bottomrule
\end{tabular}
\end{center}

\paragraph{Dữ liệu.}
Tập dữ liệu \emph{Harvard Medical Image Fusion} \cite{zhao2023cddfuse} đã được chia sẵn train/test:
\begin{itemize}\itemsep0pt
    \item \textbf{Huấn luyện}: 738 cặp ảnh ($160$ CT-MRI + $245$ PET-MRI + $333$ SPECT-MRI). Sau khi crop patch $128 \times 128$ với stride $64$ và loại bỏ patch low-contrast, còn lại $\sim 6000$--$8000$ patches dùng cho training.
    \item \textbf{Đánh giá}: 72 cặp test (24 mỗi modality CT, PET, SPECT), dùng chung với các phương pháp SOTA trong Bảng~\ref{tab:sota-raw-metrics}.
\end{itemize}

\section{Luồng xử lý dữ liệu (Data Flow)}
\label{sec:dataflow}

Phần này mô tả chi tiết cách dữ liệu đầu vào được biến đổi qua các bước xử lý, từ ảnh PNG gốc đến ảnh tổng hợp đầu ra. Luồng được chia thành hai giai đoạn: \emph{Tiền xử lý dữ liệu huấn luyện} và \emph{Suy luận trên ảnh test}.

\subsection{Luồng xử lý huấn luyện}
\label{ssec:train-dataflow}

\begin{figure}[h]
\centering
\fbox{\parbox{0.92\textwidth}{\centering\small
\texttt{[Ảnh PNG gốc]} $\to$ \texttt{[Convert sang grayscale Y]} $\to$ \texttt{[Chuẩn hóa 0-1]} \\
$\to$ \texttt{[Cắt patch 128$\times$128 stride 64]} $\to$ \texttt{[Lọc low-contrast]} $\to$ \texttt{[Lưu HDF5]} \\
$\to$ \texttt{[DataLoader batch 16]} $\to$ \texttt{[Forward train pha I/II]}
}}
\caption{Luồng tiền xử lý dữ liệu huấn luyện cho CDDFuse-AG.}
\label{fig:train-dataflow}
\end{figure}

Chi tiết từng bước:

\paragraph{Bước 1 --- Đọc ảnh raw.}
Mỗi cặp huấn luyện gồm một ảnh source ($I_V$: CT hoặc PET hoặc SPECT) và một ảnh MRI ($I_I$), kích thước $H \times W = 256 \times 256$. CT/MRI ở định dạng grayscale 8-bit; PET/SPECT ở định dạng RGB 8-bit (mã hóa màu giả --- pseudo-color).

\paragraph{Bước 2 --- Chuẩn hóa sang miền chung.}
Để đưa hai modality về cùng không gian xử lý:
\begin{itemize}\itemsep0pt
    \item Với PET/SPECT RGB: chuyển đổi không gian màu sang YCbCr, lấy kênh độ sáng $Y$:
    \begin{equation}
    Y = 0.299 R + 0.587 G + 0.114 B
    \end{equation}
    Hai kênh $C_b, C_r$ được lưu lại để khôi phục màu khi xuất ảnh.
    \item Tất cả tensor được chuẩn hóa về miền $[0, 1]$: $\tilde{I} = I / 255$.
\end{itemize}

\paragraph{Bước 3 --- Trích xuất patch.}
Mỗi cặp ảnh $(I_V, I_I)$ kích thước $256 \times 256$ được cắt thành các patch $128 \times 128$ với bước trượt (stride) bằng $64$. Số patch sinh ra trên mỗi ảnh là $\lfloor (256-128)/64 \rfloor + 1 = 3$ ngang $\times$ $3$ dọc $= 9$ patch/ảnh.

\paragraph{Bước 4 --- Lọc patch low-contrast.}
Patch có dynamic range thấp (background đồng nhất, không chứa cấu trúc) bị loại bỏ:
\begin{equation}
\text{Loại bỏ nếu} \quad \frac{p_{90} - p_{10}}{p_{90} + \epsilon} < 0.1
\end{equation}
trong đó $p_{10}, p_{90}$ là percentile $10$\% và $90$\% của patch. Mục đích: tránh huấn luyện trên patch không có thông tin, ổn định hóa decomposition loss.

\paragraph{Bước 5 --- Lưu HDF5.}
Các patch hợp lệ được lưu vào file HDF5 (\texttt{MIF\_train\_imgsize\_128\_stride\_64.h5}) với hai nhóm:
\begin{itemize}\itemsep0pt
    \item \texttt{mri\_patchs/}: tensor $[1, 128, 128]$ cho mỗi patch MRI.
    \item \texttt{src\_patchs/}: tensor $[1, 128, 128]$ cho mỗi patch source.
\end{itemize}
Định dạng này cho phép load nhanh và memory-efficient.

\paragraph{Bước 6 --- DataLoader.}
Trong vòng huấn luyện, \texttt{PyTorch DataLoader} đọc từ HDF5, tạo batch kích thước $16$ (theo paper text \S5.1), shuffle ngẫu nhiên, pin memory cho GPU transfer.

\paragraph{Bước 7 --- Forward pass.}
Tùy theo pha huấn luyện (Pha I hoặc Pha II), batch được forward qua các thành phần khác nhau (xem Mục~\ref{ssec:cddfuse-train}). Loss được tính, backward, update tham số.

\subsection{Luồng suy luận trên ảnh test}
\label{ssec:inference-dataflow}

\begin{figure}[h]
\centering
\fbox{\parbox{0.92\textwidth}{\centering\small
\texttt{[Ảnh PNG]} $\to$ \texttt{[Convert grayscale Y]} $\to$ \texttt{[Normalize 0-1]} $\to$ \texttt{[Tensor [1,1,H,W]]} \\
$\to$ \texttt{[Encoder $\to$ (f$^B$, f$^D$) cho cả 2 modality]} \\
$\to$ \texttt{[Adaptive Gating: g$^B \odot$ f$^B_V$ + (1-g$^B$) $\odot$ f$^B_I$]} (Module A) \\
$\to$ \texttt{[BaseFuseLayer, DetailFuseLayer]} \\
$\to$ \texttt{[Decoder concat $\to$ refine $\to$ sigmoid]} \\
$\to$ \texttt{[Scale 0-255 uint8]} $\to$ \texttt{[Restore màu YCbCr cho PET/SPECT]} $\to$ \texttt{[Lưu PNG]}
}}
\caption{Luồng suy luận trên ảnh test, từ đầu vào đến ảnh tổng hợp đầu ra.}
\label{fig:inference-dataflow}
\end{figure}

Chi tiết từng bước (cho một cặp ảnh test):

\paragraph{Bước 1 --- Đọc ảnh test.}
Ảnh nguồn $I_V$ và ảnh giải phẫu $I_I$ (MRI) được load từ thư mục test bằng PIL.

\paragraph{Bước 2 --- Tiền xử lý.}
\begin{itemize}\itemsep0pt
    \item Convert sang grayscale Y nếu là RGB (PET/SPECT): $Y = 0.299R + 0.587G + 0.114B$.
    \item Chuẩn hóa $[0, 1]$.
    \item Tạo tensor 4D $[1, 1, H, W]$ chuyển lên GPU.
\end{itemize}

\paragraph{Bước 3 --- Encoder phân rã.}
Hai modality được đưa qua \emph{cùng một} Restormer Encoder (shared weights):
\begin{align}
(f_V^B, f_V^D) &= \mathrm{Encoder}(I_V) \\
(f_I^B, f_I^D) &= \mathrm{Encoder}(I_I)
\end{align}
Output là hai cặp đặc trưng: Base ($64$ kênh, low-frequency, shared structure) và Detail ($64$ kênh, high-frequency, modality-specific).

\paragraph{Bước 4 --- Adaptive Gating (đóng góp của đồ án).}
Mỗi nhánh được kết hợp bằng cổng học được thay vì phép cộng đơn giản:
\begin{align}
g^B &= \sigma\big(W_g^B \cdot [f_V^B ; f_I^B] + b_g^B\big) \in (0, 1)^{64 \times H \times W} \\
\tilde{f}_F^B &= g^B \odot f_V^B + (1 - g^B) \odot f_I^B
\end{align}
Tương tự cho $\tilde{f}_F^D$. Cổng $g$ được học per-pixel và per-channel, cho phép trọng số khác nhau ở từng vị trí trong ảnh.

\paragraph{Bước 5 --- BaseFuseLayer + DetailFuseLayer.}
Đặc trưng đã gate được tinh chỉnh thêm bằng các module gốc của CDDFuse:
\begin{align}
f_F^B &= \mathrm{BaseFuseLayer}(\tilde{f}_F^B) \quad \text{(1 Transformer block)} \\
f_F^D &= \mathrm{DetailFuseLayer}(\tilde{f}_F^D) \quad \text{(1 INN block)}
\end{align}

\paragraph{Bước 6 --- Decoder tái tạo.}
Hai luồng đặc trưng được concat theo channel ($128$ kênh), giảm về $64$ kênh, refine qua $4$ Restormer block, đi qua output head và sigmoid:
\begin{equation}
\hat{I}_F^Y = \sigma\big(\mathrm{Decoder}(f_F^B, f_F^D)\big) \in [0, 1]^{H \times W}
\end{equation}

\paragraph{Bước 7 --- Hậu xử lý.}
\begin{itemize}\itemsep0pt
    \item Scale về dải $[0, 255]$ và ép kiểu \texttt{uint8}: $F^Y = \mathrm{clip}(\hat{I}_F^Y \cdot 255, 0, 255)$.
    \item Với PET/SPECT (modality màu): khôi phục màu bằng cách ghép $F^Y$ với $C_b, C_r$ của ảnh nguồn, sau đó convert YCbCr $\to$ RGB.
    \item Với CT-MRI (grayscale thuần): xuất trực tiếp $F^Y$.
\end{itemize}

\paragraph{Bước 8 --- Lưu và đánh giá.}
Ảnh tổng hợp được lưu thành file PNG. Đồng thời, 22+ chỉ số chất lượng được tính trực tiếp trên kết quả (so với $I_V, I_I$ làm reference): \texttt{NABF, SSIM, PSNR, QG, QM, MI, SCD, \ldots} (chi tiết Mục~\ref{ssec:metrics}).

\subsection{Sơ đồ tổng quan đầu-cuối}
\label{ssec:end2end-dataflow}

Hợp nhất cả huấn luyện và suy luận:

\begin{figure}[h]
\centering
\fbox{\parbox{0.95\textwidth}{\centering\small
\textbf{Huấn luyện (offline, 1 lần)}: \\
$738$ cặp ảnh train $\to$ patches $128\times 128$ $\to$ HDF5 $\to$ DataLoader (batch $16$) \\
$\to$ Train $120$ epoch (Pha I: $40$ ep decomposition $\to$ Pha II: $80$ ep fusion) \\
$\to$ Lưu checkpoint \texttt{.pth} chứa weights Encoder + Decoder + AG-Base + AG-Detail + BaseFuseLayer + DetailFuseLayer \\
\quad \\
\textbf{Suy luận (online, mỗi cặp ảnh)}: \\
Load checkpoint $\to$ đọc cặp ($I_V, I_I$) $\to$ encode $\to$ AG $\to$ refine $\to$ decode $\to$ post-process $\to$ ảnh tổng hợp.
}}
\caption{Sơ đồ end-to-end luồng dữ liệu CDDFuse-AG.}
\label{fig:end2end-dataflow}
\end{figure}

\paragraph{Baseline so sánh.}
Để đảm bảo \emph{so sánh apples-to-apples}, đồ án không chỉ so với CDDFuse pretrained gốc mà còn \emph{huấn luyện lại} CDDFuse từ đầu với cùng quy trình paper-faithful trên dữ liệu của đồ án (gọi là \textbf{CDDFuse-Paper-MIF}). So sánh chính của đồ án là:
\begin{equation*}
\boxed{\textbf{CDDFuse-AG} \quad \text{vs} \quad \textbf{CDDFuse-Paper-MIF}}
\end{equation*}
Hai mô hình dùng cùng hạ tầng, cùng dữ liệu, cùng split, cùng hyperparameters $\to$ chênh lệch chỉ phản ánh đóng góp của Adaptive Gating + Saliency-guided Pixel.

\section{Kết quả chính: CDDFuse-AG vs CDDFuse-Paper-MIF}
\label{sec:results-main}

\subsection{Số liệu trung bình theo modality}
\label{ssec:per-modal-numbers}

Bảng~\ref{tab:cddfuse-ag-per-modal} trình bày giá trị trung bình của 7 chỉ số quan trọng nhất trên từng modality. Giá trị \textbf{in đậm} là giá trị tốt hơn giữa hai mô hình.

\begin{table}[h]
\centering
\small
\caption{Giá trị trung bình theo modality (CDDFuse-AG-45ep, full 738 train pairs, $\alpha_3=10$ paper text). $\uparrow$ cao tốt hơn, $\downarrow$ thấp tốt hơn. \textbf{In đậm} = tốt hơn.}
\label{tab:cddfuse-ag-per-modal}
\setlength{\tabcolsep}{5pt}
\begin{tabular}{l|cc|cc|cc}
\toprule
& \multicolumn{2}{c|}{\textbf{CT}} & \multicolumn{2}{c|}{\textbf{PET}} & \multicolumn{2}{c}{\textbf{SPECT}} \\
\textbf{Metric} & Baseline & CDDFuse-AG & Baseline & CDDFuse-AG & Baseline & CDDFuse-AG \\
\midrule
SSIM $\uparrow$   & \textbf{1.4434} & 1.4101          & \textbf{1.3235} & 1.2924          & \textbf{1.4031} & 1.3946          \\
PSNR $\uparrow$   & \textbf{14.40}  & 14.12           & \textbf{17.53}  & 16.69           & \textbf{21.40}  & 21.20           \\
QG   $\uparrow$   & 0.6183          & \textbf{0.6454} & 0.7638          & \textbf{0.7641} & 0.7503          & \textbf{0.7514} \\
QM   $\uparrow$   & 0.0046          & \textbf{0.0148} & 0.1230          & \textbf{0.2318} & 0.0788          & \textbf{0.0909} \\
NABF $\downarrow$ & \textbf{0.0323} & 0.0356          & \textbf{0.0153} & 0.0276          & \textbf{0.0194} & 0.0202          \\
QSF  $\uparrow$   & 0.0488          & \textbf{0.0871} & \textbf{0.0546} & 0.0219          & \textbf{0.0460} & 0.0459          \\
QMI  $\uparrow$   & 0.7830          & \textbf{0.7943} & 0.8020          & \textbf{0.8094} & 0.9049          & \textbf{0.9189} \\
MI   $\uparrow$   & \textbf{53.68}  & 51.92           & 68.33           & \textbf{68.77} & \textbf{40.50}  & 40.14           \\
VAR  $\uparrow$   & \textbf{70.51}  & 70.11           & 77.51           & \textbf{78.55} & \textbf{53.82}  & 53.39           \\
\bottomrule
\end{tabular}
\end{table}

Bảng~\ref{tab:cddfuse-ag-pct} dưới đây thể hiện cùng số liệu dưới dạng \textbf{phần trăm cải thiện} so với baseline. Cell \textbf{in đậm} là cải thiện $> +3\%$ (đáng chú ý).

\begin{table}[h]
\centering
\small
\caption{Phần trăm cải thiện của CDDFuse-AG-45ep so với CDDFuse-Paper-MIF (đã hiệu chỉnh hướng cho metric ``thấp tốt hơn''). \textbf{In đậm} = cải thiện $> +3\%$; \textit{nghiêng} = regression $< -3\%$.}
\label{tab:cddfuse-ag-pct}
\begin{tabular}{lccc}
\toprule
\textbf{Metric} & \textbf{CT} & \textbf{PET} & \textbf{SPECT} \\
\midrule
SSIM   & $-2.31\%$            & $-2.35\%$            & $-0.61\%$ \\
PSNR   & $-1.93\%$            & \textit{$-4.78\%$}   & $-0.92\%$ \\
QG     & \textbf{+4.39\%}     & $+0.04\%$            & $+0.15\%$ \\
QM     & \textbf{+219.82\%}   & \textbf{+88.56\%}    & \textbf{+15.37\%} \\
NABF ($\downarrow$) & \textit{$-10.04\%$} & \textit{$-79.78\%$} & $-3.80\%$ \\
QSF    & \textbf{+78.41\%}    & \textit{$-59.91\%$}  & $-0.09\%$ \\
QMI    & $+1.45\%$            & $+0.92\%$            & $+1.55\%$ \\
MI     & $-3.27\%$            & $+0.65\%$            & $-0.89\%$ \\
VAR    & $-0.55\%$            & $+1.34\%$            & $-0.80\%$ \\
EN     & $-1.94\%$            & $-2.52\%$            & $-2.58\%$ \\
\bottomrule
\end{tabular}
\end{table}

\paragraph{Quan sát chính.}
\begin{itemize}\itemsep0pt
    \item \textbf{CT}: cải thiện rõ rệt ở \emph{QM $+219\%$} (wavelet quality tăng rất mạnh), \emph{QSF $+78\%$} (sharpness), \emph{QG $+4.4\%$} (gradient). Đánh đổi nhẹ ở SSIM/PSNR ($\sim -2\%$) và NABF tăng nhẹ.
    \item \textbf{PET}: \emph{QM $+89\%$} (wavelet rất mạnh), QMI/QG tăng nhẹ. Nhưng QSF $-60\%$ và NABF $-80\%$ (artifact tăng) — trade-off rõ.
    \item \textbf{SPECT}: \emph{QM $+15\%$}, QMI $+1.55\%$ cải thiện ổn định. Các metric khác gần như tương đương ($\pm 1\%$).
\end{itemize}

\subsection{Trung bình tổng quan trên toàn tập (Pooled)}
\label{ssec:pooled-results}

Khi gộp toàn bộ 72 cặp test (24 mỗi modality), giá trị trung bình \emph{toàn tập} của hai mô hình:

\begin{table}[h]
\centering
\small
\caption{Trung bình trên 72 cặp test gộp 3 modality (CDDFuse-AG-45ep, full 738 train, $\alpha_3=10$). \textbf{In đậm} = tốt hơn.}
\label{tab:cddfuse-ag-pooled}
\begin{tabular}{lccc}
\toprule
\textbf{Metric} & \textbf{Baseline} & \textbf{CDDFuse-AG} & \textbf{Cải thiện (\%)} \\
\midrule
SSIM $\uparrow$   & \textbf{1.3900} & 1.3657          & $-1.75\%$ \\
PSNR $\uparrow$   & \textbf{17.78}  & 17.34           & $-2.46\%$ \\
QG   $\uparrow$   & 0.7108          & \textbf{0.7203} & $+1.34\%$ \\
QM   $\uparrow$   & 0.0688          & \textbf{0.1125} & \textbf{$+63.56\%$} \\
NABF $\downarrow$ & \textbf{0.0224} & 0.0278          & $-24.19\%$ (artifact tăng) \\
QSF  $\uparrow$   & 0.0498          & \textbf{0.0516} & $+3.68\%$ \\
QMI  $\uparrow$   & 0.8300          & \textbf{0.8409} & $+1.31\%$ \\
MI   $\uparrow$   & \textbf{54.17}  & 53.61           & $-1.03\%$ \\
EN   $\uparrow$   & \textbf{5.04}   & 4.92            & $-2.36\%$ \\
VAR  $\uparrow$   & 67.28           & \textbf{67.35}  & $+0.11\%$ ($\approx$ tie) \\
\bottomrule
\end{tabular}
\end{table}

\paragraph{Đánh giá tổng quan.}
Trên 72 cặp test gộp 3 modality (CDDFuse-AG-45ep):
\begin{itemize}\itemsep0pt
    \item \textbf{QM tăng $63.6\%$} (wavelet quality) --- đây là điểm cải thiện mạnh nhất, cho thấy AG preserve chi tiết tần số cao tốt hơn baseline.
    \item \textbf{QSF tăng $3.7\%$, QG tăng $1.3\%$, QMI tăng $1.3\%$}: cấu trúc gradient, sharpness và mutual information đều nhỉnh hơn.
    \item Trade-off: NABF tăng $24.2\%$ (artifact nhẹ), SSIM/PSNR giảm $\sim 2\%$, EN giảm $2.4\%$.
\end{itemize}
Đây là trade-off đặc trưng với $\alpha_3=10$ (TV gradient loss mạnh): mô hình ưu tiên giữ \emph{chi tiết wavelet} và \emph{sharpness} (QM, QSF, QG cải thiện) thay vì làm sạch artifact biên (NABF). Trade-off này phù hợp khi mục tiêu lâm sàng là \emph{quan sát rõ texture vi mô} (tổn thương, vùng nghi ngờ).

\subsection{Kiểm định ý nghĩa thống kê}
\label{ssec:stats-results}

Áp dụng pipeline Wilcoxon + Cliff's $\delta$ + Holm--Bonferroni (xem Mục~\ref{sec:stats-bg}) trên 25 chỉ số chất lượng, kết quả \emph{số metric có ý nghĩa thống kê} (SIG):

\begin{table}[h]
\centering
\small
\caption{Số chỉ số có ý nghĩa thống kê sau Holm correction ($\alpha = 0.05$, $K = 22$). Per-modality test riêng từng modality ($n = 24$); Pooled gộp 3 modality ($n = 72$).}
\label{tab:cddfuse-ag-sig}
\begin{tabular}{lcccc}
\toprule
& \textbf{CT} & \textbf{PET} & \textbf{SPECT} & \textbf{Pooled} \\
\midrule
Số SIG / 22 metric & \textbf{8} + 1 MARG & \textbf{6} + 1 MARG & \textbf{5} + 6 MARG & \textbf{8} \\
Verdict           & CONFIRM    & CONFIRM   & CONFIRM & \textbf{CONFIRM} \\
\bottomrule
\end{tabular}
\end{table}

\paragraph{Diễn giải.}
\begin{itemize}\itemsep0pt
    \item \textbf{Cả 3 modality đạt CONFIRM\_IMPROVEMENT} ($\ge 5$ SIG ngưỡng): CDDFuse-AG cải thiện ý nghĩa thống kê ở \emph{cả CT, PET và SPECT}, đây là kết quả mạnh nhất.
    \item \textbf{Pooled 8 SIG (CONFIRM)}: hiệu quả \emph{nhất quán} trên 72 cặp gộp, không phải gain riêng modality nào.
    \item Top metric SIG (pooled): \textbf{QM} ($\delta = +0.285$ small), \textbf{CE} ($+0.169$), \textbf{QSF} ($+0.173$), \textbf{SF} ($+0.084$), \textbf{QMI} ($+0.079$), \textbf{AG} ($+0.076$), \textbf{QG} ($+0.118$).
    \item Top regression (NS sau Holm, $|\delta| > 0.2$): QY, NABF, SSIM --- liên quan đến perceptual quality và artifact.
\end{itemize}

\section{Phân tích trade-off}
\label{sec:results-tradeoff}

Quan sát từ kết quả CDDFuse-AG-45ep (full 738 train, $\alpha_3=10$): mô hình mạnh ở chỉ số \emph{wavelet + frequency + gradient} (QM, QSF, SF, QG) nhưng giảm nhẹ ở chỉ số \emph{perceptual + structure} (SSIM, PSNR, NABF). Điều này hợp lý với cơ chế:
\begin{itemize}\itemsep0pt
    \item Hệ số TV loss $\alpha_3 = 10$ (paper text, lớn hơn 5 trong code release) ép mô hình \emph{giữ gradient/edges/details} $\to$ QM tăng $63\%$, QSF/QG tăng nhẹ.
    \item Adaptive Gating chọn lọc kênh đặc trưng $\to$ tăng chất lượng wavelet và mutual information.
    \item Saliency-guided Pixel target ưu tiên vùng có gradient cao $\to$ output sắc nét ở biên, nhưng tạo \emph{ridge mỏng} ở một số vùng làm NABF nhỉnh tăng và SSIM giảm nhẹ.
\end{itemize}
Đây là trade-off đặc trưng của mô hình ưu tiên \emph{texture vi mô và edge preservation} (mục tiêu chẩn đoán: nhìn rõ tổn thương) thay vì \emph{smoothing perceptual} (giảm artifact đồng đều).

\section{So sánh phụ với CDDFuse pretrained công bố}
\label{sec:results-vs-pretrained}

Để đặt CDDFuse-AG trong bối cảnh SOTA, đồ án cũng so sánh với checkpoint pretrained CDDFuse công bố (\texttt{CDDFuse\_MIF.pth}, không retrain). Đây không phải so sánh apples-to-apples (khác training data composition) nên chỉ là tham khảo.

\begin{table}[h]
\centering
\small
\caption{So với CDDFuse pretrained công bố --- tham khảo. \textbf{In đậm} = CDDFuse-AG tốt hơn pretrained.}
\label{tab:vs-pretrained}
\begin{tabular}{lccc}
\toprule
\textbf{Modal} & \textbf{Cải thiện nổi bật} & \textbf{Verdict (Wilcoxon)} & \textbf{Trade-off} \\
\midrule
CT    & NABF giảm $\sim 8\%$, QM tăng $\sim 8\%$, QMI tăng $\sim 2\%$ & 9 SIG / 22 (CONFIRM) & QSF giảm $\sim 23\%$, MI giảm $\sim 4\%$ \\
PET   & NABF giảm $\sim 7\%$, QSF tăng $\sim 2\%$ & 11 SIG / 22 (CONFIRM) & QM giảm $\sim 9\%$, MI giảm $\sim 1\%$ \\
SPECT & tương đương baseline ($\pm 2\%$) & 2 SIG / 22 (MIXED) & QSF giảm $\sim 4\%$ \\
\bottomrule
\end{tabular}
\end{table}

So với SOTA Bảng~\ref{tab:sota-raw-metrics}, CDDFuse-AG củng cố thêm vị trí top tier của CDDFuse, đặc biệt ở các chỉ số mà CDDFuse gốc còn yếu (NABF, QSF tùy modality).
